\chapter{Fundamentação Teórica}
\label{cap:fundamentacao}

Este capítulo apresenta os fundamentos conceituais necessários à compreensão da proposta
deste trabalho, organizados da seguinte forma: as Redes Ad Hoc e sua especialização em
Redes Veiculares Ad Hoc (VANETs); a arquitetura e os modelos de mobilidade característicos
desse tipo de rede; os protocolos de roteamento empregados, com ênfase no roteamento
geográfico; e, por fim, o simulador ns-3, sua arquitetura, sua extensibilidade e o módulo
FlowMonitor, utilizado para coleta de métricas de desempenho.

\section{Redes Ad Hoc}
\label{sec:redes-ad-hoc}

Uma rede ad hoc é uma rede sem fio autônoma, formada dinamicamente por um conjunto de
dispositivos móveis que se comunicam diretamente entre si, sem depender de uma
infraestrutura de rede fixa ou de um ponto de acesso centralizado. Cada nó atua
simultaneamente como terminal e como roteador, encaminhando pacotes em nome de outros nós
quando estes não estão ao alcance direto do destinatário. Essa característica confere às
redes ad hoc grande flexibilidade de implantação, mas também impõe desafios específicos de
roteamento, uma vez que a topologia da rede pode se alterar constantemente à medida que os
nós se movimentam.

As Redes Móveis Ad Hoc (MANETs, do inglês \textit{Mobile Ad Hoc Networks}) constituem a
categoria mais geral desse tipo de rede, sendo compostas por dispositivos móveis diversos
(notebooks, sensores, smartphones) sem restrição quanto ao padrão de mobilidade.
\citeonline{pimenta2013simulacao} avalia a simulação de diversos protocolos de roteamento
propostos para MANETs, evidenciando que o desempenho desses protocolos está diretamente
relacionado à intensidade e ao padrão de mobilidade dos nós da rede — aspecto que se torna
ainda mais crítico no subtipo de rede tratado neste trabalho, as VANETs, descritas na seção
seguinte.

\section{Redes Veiculares Ad Hoc (VANETs)}
\label{sec:vanets}

As Redes Veiculares Ad Hoc (VANETs) constituem um subtipo especializado de MANET, no qual
os nós móveis são veículos equipados com dispositivos de comunicação sem fio. Diferem das
MANETs genéricas por apresentarem padrões de mobilidade restritos por vias, velocidades
relativas elevadas entre os nós e uma topologia de rede altamente dinâmica, decorrente das
frequentes rupturas de enlace causadas pelo deslocamento dos veículos.

\citeonline{meneguette2013seamless} discutem os desafios de gerenciamento de mobilidade em
redes de comunicação veicular, evidenciando que arquiteturas de rede tradicionais
frequentemente não lidam de forma satisfatória com as transições de conectividade
motivadas pela alta velocidade dos nós. As VANETs se inserem no contexto mais amplo dos
Sistemas de Transporte Inteligentes (ITS, do inglês \textit{Intelligent Transportation
Systems}), sendo aplicadas em cenários como prevenção de colisões, disseminação de alertas
de trânsito e otimização de fluxo viário.

\section{Arquitetura de uma VANET}
\label{sec:arquitetura-vanet}

A arquitetura de uma VANET é tipicamente composta por três elementos principais: a Unidade
de Bordo (\textit{On-Board Unit}, OBU), instalada em cada veículo e responsável pela
comunicação sem fio e pelo processamento das mensagens trocadas; a Unidade de Acostamento
(\textit{Roadside Unit}, RSU), instalada em pontos fixos ao longo das vias e responsável por
estender a cobertura da rede e por conectá-la à infraestrutura de backend; e, eventualmente,
uma Autoridade Confiável, responsável por aspectos de segurança e gerenciamento de
identidade dos veículos.

Essa arquitetura sustenta dois domínios de comunicação complementares, já mencionados na
Seção~\ref{sec:vanets}: a comunicação veículo-veículo (V2V), estabelecida diretamente entre
OBUs sem depender de infraestrutura fixa; e a comunicação veículo-infraestrutura (V2I),
estabelecida entre OBUs e RSUs. A arquitetura de mobilidade de fluxo discutida por
\citeonline{meneguette2013seamless} ilustra como a transição entre esses dois domínios de
comunicação — por exemplo, quando um veículo se afasta de uma RSU e passa a depender
exclusivamente de enlaces V2V — constitui um dos principais desafios de projeto em VANETs,
com impacto direto sobre a escolha do protocolo de roteamento empregado.

\section{Mobilidade em VANETs}
\label{sec:mobilidade}

A mobilidade dos nós é o fator que mais distingue as VANETs de outras redes ad hoc, e sua
correta modelagem é essencial para a validade de qualquer avaliação experimental. Duas
abordagens principais são identificadas na literatura para a simulação de mobilidade
veicular no ns-3.

A primeira consiste no uso de modelos de mobilidade nativos do próprio simulador.
\citeonline{arbabi2010highway} apresentam um modelo de mobilidade voltado especificamente à
simulação de rodovias, integrando mobilidade e comunicação de rede por meio do mecanismo de
eventos discretos do ns-3, sem depender de ferramentas externas.

A segunda abordagem, mais frequente na literatura recente, integra o ns-3 a um simulador de
mobilidade dedicado, como o SUMO (\textit{Simulation of Urban Mobility}), por meio do
protocolo TraCI. Essa integração é adotada por \citeonline{kolici2015performance},
\citeonline{su2014improved} e \citeonline{muktar2025machine}, permitindo reproduzir cenários
de tráfego urbano mais realistas, com veículos respondendo a semáforos, congestionamentos e
regras de trânsito, à custa de maior complexidade de configuração experimental.
\citeonline{sharma2026modeling} revisam quatorze modelos de mobilidade distintos aplicados à
avaliação de protocolos de roteamento em VANETs, reforçando que a escolha do modelo de
mobilidade influencia diretamente os resultados de desempenho obtidos. Neste trabalho, opta-se
pelo uso de modelos de mobilidade nativos do ns-3, pelas razões metodológicas apresentadas no
Capítulo~\ref{cap:metodologia}.

\section{Protocolos de Roteamento em VANETs}
\label{sec:protocolos-roteamento}

Os protocolos de roteamento em redes ad hoc — e, por extensão, em VANETs — podem ser
classificados em duas categorias principais: protocolos topológicos, que mantêm tabelas de
rotas construídas a partir da topologia da rede; e protocolos geográficos, que utilizam a
posição física dos nós como critério para o encaminhamento de pacotes, tema tratado em
detalhe na Seção~\ref{sec:roteamento-geografico}.

Entre os protocolos topológicos, destacam-se o AODV (\textit{Ad hoc On-Demand Distance
Vector}) e o OLSR (\textit{Optimized Link State Routing}), ambos nativamente suportados pelo
ns-3 e amplamente utilizados como referência de comparação em estudos de desempenho de redes
ad hoc. \citeonline{pimenta2013simulacao} demonstra que, embora sejam soluções consolidadas,
esses protocolos tendem a apresentar desempenho degradado em cenários de alta mobilidade, uma
vez que a reconstrução das tabelas de rotas após rupturas de enlace consome tempo e largura
de banda adicionais — limitação especialmente relevante no contexto das VANETs, dada a
mobilidade intensa discutida na Seção~\ref{sec:mobilidade}.

\section{Roteamento Geográfico}
\label{sec:roteamento-geografico}

Como alternativa aos protocolos topológicos, os protocolos de roteamento geográfico
prescindem da manutenção de tabelas de rotas completas, utilizando apenas a posição física
dos nós — tipicamente obtida via GPS — para decidir o encaminhamento de cada pacote. Essa
característica os torna particularmente adequados a cenários de alta mobilidade, nos quais a
manutenção de informação topológica atualizada tem custo elevado.

Dentre os protocolos geográficos, o GPSR (\textit{Greedy Perimeter Stateless Routing}) é um
dos mais estudados na literatura, operando em dois modos complementares: o modo guloso
(\textit{greedy forwarding}), em que o pacote é encaminhado ao vizinho geograficamente mais
próximo do destino; e o modo de recuperação por perímetro (\textit{perimeter/face routing}),
acionado quando o modo guloso encontra um mínimo local, isto é, quando não há vizinho mais
próximo do destino do que o próprio nó atual. \citeonline{katsaros_positionbased_ns3}
documentaram o desenvolvimento de um módulo de roteamento baseado em posição para o ns-3,
estabelecendo a arquitetura de referência do GPSR no simulador — implementação que constitui
o objeto central da modernização proposta neste trabalho.

O interesse contínuo da comunidade científica em aprimorar o GPSR resultou em diversas
variantes do protocolo original, entre as quais se destacam: o AGPSR
(\textit{Adaptive GPSR}), proposto por \citeonline{silva2018adaptive}, que introduz seleção
adaptativa de próximo salto para reduzir a taxa de pacotes perdidos em cenários de alta
densidade veicular; o PA-GPSR (\textit{Path-Aware GPSR}), de \citeonline{zhang2023weight}, que
pondera a escolha de rotas com base em métricas adicionais de qualidade de enlace; o WA-GPSR
(\textit{Weight-Aware GPSR}), de \citeonline{smiri2021wagpsr}, que introduz um esquema de
pesos para equilibrar múltiplos critérios de encaminhamento; e o VANET-GPSR+, proposto por
\citeonline{zhang2026vanet}, uma variante recente que incorpora ciência de direção
(\textit{direction-aware}) ao processo de decisão. Essas variantes evidenciam que o GPSR
permanece uma linha de pesquisa ativa, reforçando a relevância de disponibilizar à comunidade
uma implementação de referência atualizada e funcional do algoritmo original no ns-3.

\section{O Simulador ns-3}
\label{sec:ns3}

O ns-3 (\textit{Network Simulator 3}) é uma ferramenta de simulação de eventos discretos
voltada à pesquisa e ao ensino em redes de computadores, amplamente utilizada tanto na
indústria quanto na academia por ser um software livre e de código aberto.

\subsection{Arquitetura do ns-3}
\label{sec:arquitetura-ns3}

O simulador é estruturado em módulos escritos em C++, cada um correspondente a uma
funcionalidade de rede específica — por exemplo, uma pilha de protocolos, um tipo de
dispositivo de rede ou um modelo de mobilidade. Esses módulos podem ser habilitados,
desabilitados ou combinados conforme as necessidades do cenário experimental, que é descrito
por meio de scripts em C++ ou em Python que instanciam nós, atribuem-lhes módulos de rede e
mobilidade, e configuram os eventos da simulação.

\subsection{Extensibilidade do ns-3}
\label{sec:extensibilidade-ns3}

Além do uso dos módulos nativos, o ns-3 permite que novos módulos sejam desenvolvidos e
integrados ao simulador, estendendo suas funcionalidades. O estudo sistemático da literatura
conduzido por \citeonline{campanile2020computer} mapeou o uso do ns-3 em 128 publicações
científicas, constatando que essa extensibilidade é praticada por apenas uma minoria dos
trabalhos analisados — a maior parte se limita ao uso dos módulos já disponíveis nativamente.
Iniciativas voltadas ao ensino, como o portal de aplicação proposto por
\citeonline{maciel2014proposta}, e estudos comparativos entre simuladores, como o de
\citeonline{silva2017avaliacao}, reforçam esse cenário no contexto brasileiro. De forma
semelhante, \citeonline{conterato2014avaliacao} relatam limitações na avaliação do suporte do
ns-3 a redes definidas por software baseadas em OpenFlow, e \citeonline{fernandes2012scalable}
discutem restrições de escalabilidade do simulador especificamente em cenários de VANETs de
grande porte — evidências de que funcionalidades não incorporadas ao núcleo oficial do
simulador tendem a ficar defasadas ou dependentes de manutenção não sistemática pela
comunidade.

\subsection{APIs Relacionadas ao Roteamento}
\label{sec:apis-roteamento}

A criação de um novo protocolo de roteamento no ns-3 exige a implementação da interface
\texttt{Ipv4RoutingProtocol}, que define os métodos responsáveis pelo encaminhamento
(\texttt{RouteOutput}) e pelo recebimento (\texttt{RouteInput}) de pacotes IP, além de
métodos de notificação de mudanças na interface de rede. Essa API, assim como o sistema de
\textit{build} do simulador — que migrou do \textit{waf}, utilizado em versões mais antigas,
para o CMake, adotado nas versões atuais —, sofreu alterações ao longo das diferentes versões
do ns-3. A implementação do GPSR documentada por
\citeonline{katsaros_positionbased_ns3} foi concebida sob a API e o sistema de \textit{build}
de uma versão antiga do simulador, sendo, por esse motivo, incompatível com as versões atuais
— incompatibilidade que motiva o processo de modernização detalhado no
Capítulo~\ref{cap:implementacao}.

\section{FlowMonitor}
\label{sec:flowmonitor}

Para a coleta de métricas de desempenho de rede no ns-3, utiliza-se o módulo FlowMonitor,
desenvolvido por \citeonline{carneiro2009flowmonitor}, que permite a detecção automática de
fluxos de tráfego e o registro persistente de métricas como taxa de transferência
(\textit{throughput}), atraso médio (\textit{delay}), variação do atraso (\textit{jitter}) e
taxa de perda de pacotes. O módulo opera de forma não intrusiva, monitorando os pacotes que
atravessam os nós da simulação sem exigir modificações nos protocolos avaliados, o que o
torna adequado à comparação entre diferentes algoritmos de roteamento sob as mesmas condições
experimentais — característica explorada na avaliação comparativa descrita no
Capítulo~\ref{cap:metodologia}.